\documentclass{article}
\usepackage[T1]{fontenc}
\usepackage{lmodern}
\usepackage{graphicx}
\usepackage{amsmath,amssymb}
\usepackage[a4paper,margin=2cm]{geometry}
\usepackage[absolute,overlay]{textpos}
\setlength{\TPHorizModule}{1mm}
\setlength{\TPVertModule}{1mm}
\usepackage[indonesian]{babel}
\usepackage{hyperref}
\setlength{\emergencystretch}{2em}
% unit: Halaman 21

\begin{document}

% === Halaman 21 (llm) ===

Sebuah algoritma kriptografi dikatakan aman secara komputasi (\textit{computationally secure}) bila ia memenuhi tiga kriteria berikut:

\begin{enumerate}
    \item Persamaan matematika yang menggambarkan operasi di dalam algoritma kriptografi sangat kompleks sehingga algoritma tidak mungkin dipecahkan secara analitik.
    \item Biaya untuk memecahkan cipherteks melampaui nilai informasi yang terkandung di dalam cipherteks tersebut.
    \item Waktu yang diperlukan untuk memecahkan cipherteks melampaui lamanya waktu informasi tersebut harus dijaga kerahasiaannya.
\end{enumerate}

\end{document}
